\section{Conclusion}\label{sec:conclusion}

Holding the preprocessing, the architecture and the training protocol fixed and varying
only the provenance of the label shows that the accuracy ordinarily reported for cough
detection is a property of the evaluation as much as of the detector. A ResNet18 that
recovers $97.72$ per cent of coughs from held out speakers in a prompted corpus recovers
$56.04$ per cent of them from human annotated audio at the same class prior, with no
retraining between the two measurements, while its specificity falls far less; the
detector does not stop working so much as stop finding coughs, and a single accuracy
figure hides which of the two has happened. Wavelet shrinkage does not close the gap, and
the reason is measurable rather than conjectural: on recordings of this quality the
universal threshold removes a median of $0.67$ per cent of signal power, so the two input
conditions it produces are very nearly the same images and no difference between them was
available to be found. The classifier derived regime cannot settle the question either,
because its published selection is not reproducible by any threshold on the score it
supplies and its expert annotated subset is drawn overwhelmingly from the region where
that score was already confident, which leaves the label error rate of the corpus
unestimable from what it publishes. None of this says that cough detection is
unachievable, and the networks examined here fit the prompted corpus almost perfectly from
a speaker disjoint partition. It says that the experiment which would have revealed the
limitation costs a few hours of inference on a corpus nobody trained on, that it was
available throughout, and that reporting the number obtained without it describes the
corpus rather than the task.
